\begin{abstract}
\textbf{Context:} AI agents are rapidly adopted in software development, but their token costs are rising and it is unclear how much a team should spend on agents, and on which tasks.\\
\textbf{Objective:} We determine how teams should divide spending between human developers and AI agents to deliver the features a project demands at the lowest total cost.\\
\textbf{Method:} We model mixed teams of human developers and AI agents with eight specialized roles as a discrete-time stochastic process covering bug injection and detection, technical debt, human supervision of agents, and release cycles with feature freezes. Using Monte Carlo simulation and evolutionary search, we determine cost-optimal teams for 432 scenarios varying developer capability, project velocity, debt dynamics, agent capability and price, and bug costs, and compare them with a human-only baseline.\\
\textbf{Results:} Escaped bugs are the dominant cost driver, and 61\% of the agent budget goes to finding bugs. Agent capability matters more than token price, since higher-quality outputs reduce downstream quality assurance costs. In teams with capable agents, human developers mostly fix bugs and leave the implementation of new features to the agents. Teams with high-capability agents are typically cheaper than human-only teams, but low-capability agents offer no consistent advantage. At realistic velocities, the optimal teams still spend more on developers than on agents.\\
\textbf{Conclusions:} There is neither a universal token budget per developer nor a universal agentic team structure. As in traditional development, this depends on many factors and teams need to be carefully tailored to the project context to achieve the optimal cost-quality trade-off. 
\end{abstract}

\keywords{AI agents \and software process simulation \and simulation-based optimization
\and technical debt \and cost of software quality \and team composition}
